% SEGMENT OLP-0105-S01
% Part: first-order-logic
% Chapter: tableaux
% Section: proof-theoretic-notions

\documentclass[../../../include/open-logic-section]{subfiles}

\begin{document}

\iftag{FOL}
      {\olfileid{fol}{tab}{ptn}}
      {\olfileid{pl}{tab}{ptn}}
      
\olsection{நிறுவல்-கோட்பாட்டுக் கருத்துகள்}

\begin{editorial}
இந்தப் பிரிவு அட்டவணை மரங்களுக்கான நிறுவத்தக்கதன்மை உறவு மற்றும் முரணின்மை
ஆகியவற்றின் வரையறைகளைத் தொகுக்கிறது.
\end{editorial}

\begin{explain}
பல முக்கியமான பொருண்மைக் கருத்துகளை (செல்லுபடித்தன்மை, பின்விளைவு,
நிறைவுறத்தக்கதன்மை) நாம் வரையறுத்ததைப் போலவே, அவற்றுடன் தொடர்புடைய
\emph{நிறுவல்-கோட்பாட்டுக் கருத்துகளை} இப்போது வரையறுக்கிறோம். இவை
!!{structure}s[loc] !!{sentence}s நிறைவுறுவதைச் சார்ந்து வரையறுக்கப்படாமல்,
குறிப்பிட்ட மூடிய !!{tableau}s இருப்பதைச் சார்ந்து வரையறுக்கப்படுகின்றன.
இந்தக் கருத்துகள் ஒன்றுபடுகின்றன என்பது ஒரு முக்கியமான கண்டுபிடிப்பு. அவை
ஒன்றுபடுவதே \emph{சரித்தன்மைத் தேற்றம்} மற்றும் \emph{நிறைவுத் தேற்றம்}
ஆகியவற்றின் உள்ளடக்கம்.
\end{explain}


% SEGMENT OLP-0105-S02
\begin{defn}[தேற்றங்கள்]
$\sFmla{\False}{!A}$ க்கான மூடிய !!a{tableau} இருந்தால் !!{sentence}~$!A$
ஒரு \emph{தேற்றம்}. $!A$ தேற்றமாக இருந்தால் $\Proves !A$ என்றும்,
இல்லையெனில் $\Proves/ !A$ என்றும் எழுதுகிறோம்.
\end{defn}

% SEGMENT OLP-0105-S03
\begin{defn}[!!^{derivability}]
!!{sentence}s கொண்ட~$\Gamma$ கணத்திலிருந்து !!^a{sentence}~$!A$ ஆனது
\emph{!!{derivable}} என்பதை $\Gamma \Proves !A$ என்று எழுதுகிறோம்; இதன்
பொருள் $\{!B_1, \dots, !B_n\} \subseteq \Gamma$ என்ற முடிவுறு கணமும்,
பின்வரும் கணத்திற்கான மூடிய !!a{tableau} ஒன்றும் உள்ளன என்பதாகும்:
\[
\{
  \sFmla{\False}{!A},
  \sFmla{\True}{!B_1}, \dots,
  \sFmla{\True}{!B_n}
  \}.
\]
$!A$ ஆனது~$\Gamma$ இலிருந்து !!{derivable} அல்ல எனில் $\Gamma \Proves/
!A$ என்று எழுதுகிறோம்.
\end{defn}

% SEGMENT OLP-0105-S04
\begin{defn}[முரணின்மை]
!!{sentence}s கொண்ட~$\Gamma$ கணம், $\{!B_1, \dots, !B_n\} \subseteq
\Gamma$ என்ற முடிவுறு கணமும் பின்வரும் கணத்திற்கான மூடிய !!a{tableau} ஒன்றும்
இருந்தால், அப்போதும் அப்போது மட்டுமே \emph{முரணுடையது}:
\[
\{
  \sFmla{\True}{!B_1}, \dots,
  \sFmla{\True}{!B_n}  
  \}.
\]
$\Gamma$ முரணுடையதல்ல எனில், அது \emph{முரணற்றது} என்கிறோம்.
\end{defn}

% SEGMENT OLP-0105-S05
\begin{prop}[தற்சுட்டுப் பண்பு]
\ollabel{prop:reflexivity}
$!A \in \Gamma$ எனில், $\Gamma \Proves !A$.
\end{prop}

\begin{proof}
$!A \in \Gamma$ எனில், $\{!A\}$ ஆனது~$\Gamma$ வின் முடிவுறு உட்கணம்;
பின்வரும் !!{tableau} மூடியது:
\begin{oltableau}
  [\sFmla{\False}{\formula{A}}, just = \TAss
    [\sFmla{\True}{\formula{A}}, just = \TAss,close]
  ]
\end{oltableau}
\end{proof}
  

% SEGMENT OLP-0105-S06
\begin{prop}[ஒருபோக்குத்தன்மை]
\ollabel{prop:monotonicity}
$\Gamma \subseteq \Delta$ மற்றும் $\Gamma \Proves !A$ எனில், $\Delta
\Proves !A$.
\end{prop}

\begin{proof}
$\Gamma$ வின் எந்த முடிவுறு உட்கணமும்~$\Delta$ வின் முடிவுறு உட்கணமாகும்.
\end{proof}

% SEGMENT OLP-0105-S07
\begin{prop}[கடப்புப் பண்பு]
\ollabel{prop:transitivity}
$\Gamma \Proves !A$ மற்றும் $\{!A\} \cup
\Delta \Proves !B$ எனில், $\Gamma \cup \Delta \Proves !B$.
\end{prop}

\begin{proof}
$\{!A\} \cup \Delta \Proves !B$ எனில், $\Delta_0 =
\{!C_1, \dots, !C_n\} \subseteq \Delta$ என்ற முடிவுறு உட்கணம் இருந்து,
\begin{align*}
\{\sFmla{\False}{!B}, & \sFmla{\True}{!A}, \sFmla{\True}{!C_1},
\dots, \sFmla{\True}{!C_n}\}
\intertext{என்பதற்கு மூடிய !!a{tableau} உள்ளது. $\Gamma \Proves !A$ எனில்,
  பின்வருவதற்கு மூடிய !!a{tableau} இருக்குமாறு $!D_1$, \dots,
  $!D_m \subseteq \Gamma$ உள்ளன:}
\{\sFmla{\False}{!A}, & \sFmla{\True}{!D_1},
\dots, \sFmla{\True}{!D_m}\}
\end{align*}

இப்போது பின்வரும் எடுகோள்களையுடைய !!{tableau}வைக் கருதுக:
\[
\sFmla{\False}{!B},
\sFmla{\True}{!C_1}, \dots, \sFmla{\True}{!C_n},
\sFmla{\True}{!D_1}, \dots, \sFmla{\True}{!D_m}.
\]
$!A$ மீது \Cut{} விதியைப் பயன்படுத்துக. இது இரு கிளைகளை உருவாக்குகிறது;
ஒன்றில் $\sFmla{\True}{!A}$ உம், மற்றதில் $\sFmla{\False}{!A}$ உம் உள்ளன.
எனவே ஒரு கிளையில் பின்வரும் அனைத்தும் கிடைக்கின்றன:
\[
\{\sFmla{\False}{!B}, \sFmla{\True}{!A},
\sFmla{\True}{!C_1}, \dots, \sFmla{\True}{!C_n}\}
\]
இந்த எடுகோள்களுக்கு மூடிய !!a{tableau} இருப்பதால், அதை அந்தக் கிளையில்
இணைக்கலாம்; $\sFmla{\True}{!A}$ வழியாகச் செல்லும் ஒவ்வொரு கிளையும் மூடுகிறது.
மற்ற கிளையில் பின்வரும் அனைத்தும் கிடைக்கின்றன:
\[
\{\sFmla{\False}{!A}, \sFmla{\True}{!D_1}, \dots,
\sFmla{\True}{!D_m}\}
\]
எனவே மறுபக்கத்தையும் நிறைவுசெய்து மூடிய !!a{tableau}வைப் பெறலாம். இது
$\Gamma \cup \Delta \Proves !B$ என்பதைக் காட்டுகிறது.
\end{proof}

% SEGMENT OLP-0105-S08
குறிப்பாக $\Gamma \Proves !A$ மற்றும் $!A \Proves !B$ எனில், $\Gamma
\Proves !B$ என்பது இதன் பொருள். மேலும் $!A_1, \dots, !A_n \Proves !B$
என்றும், ஒவ்வொரு~$i$ க்கும் $\Gamma \Proves !A_i$ என்றும் இருந்தால்,
$\Gamma \Proves !B$ என்பதும் பின்வருகிறது.

% SEGMENT OLP-0105-S09
\begin{prop}
\ollabel{prop:incons}
$\Gamma$ முரணுடையது என்பது, அப்போதும் அப்போது மட்டுமே, ஒவ்வொரு
!!{sentence}~$!A$ க்கும் $\Gamma \Proves !A$.
\end{prop}

\begin{proof}
பயிற்சி.
\end{proof}

\tagprob{FOL}
\begin{prob}
\olref[fol][tab][ptn]{prop:incons} ஐ நிறுவுக.
\end{prob}
\tagendprob

\tagprob{notFOL}
\begin{prob}
\olref[pl][tab][ptn]{prop:incons} ஐ நிறுவுக.
\end{prob}
\tagendprob

% SEGMENT OLP-0105-S10
\begin{prop}[இறுக்கத்தன்மை]
  \ollabel{prop:proves-compact}
  \begin{enumerate}
  \item $\Gamma \Proves !A$ எனில், $\Gamma_0 \subseteq \Gamma$ என்ற ஒரு
    முடிவுறு உட்கணம் இருந்து, $\Gamma_0 \Proves !A$.
  \item $\Gamma$ வின் ஒவ்வொரு முடிவுறு உட்கணமும் முரணற்றது எனில்,
    $\Gamma$ முரணற்றது.
  \end{enumerate}
\end{prop}

\begin{proof}
  \begin{enumerate}
    \item $\Gamma \Proves !A$ எனில், $\Gamma_0 = \{!B_1, \dots, !B_n\}$
      என்ற முடிவுறு உட்கணமும் பின்வருவதற்கான மூடிய !!a{tableau}வும் உள்ளன:
      \[
      \{\sFmla{\False}{!A}, \sFmla{\True}{!B_1}, \dots, \sFmla{\True}{!B_n}\}
      \]
      இந்த !!{tableau} $\Gamma_0 \Proves !A$ என்பதையும் காட்டுகிறது.
    \item $\Gamma$ முரணுடையது எனில், $\Gamma_0 = \{!B_1, \dots, !B_n\}$
      என்ற ஏதோவொரு முடிவுறு உட்கணத்திற்குப் பின்வருவதற்கான மூடிய
      !!a{tableau} உள்ளது:
      \[
      \{\sFmla{\True}{!B_1}, \dots, \sFmla{\True}{!B_n}\}
      \]
      இந்த மூடிய !!{tableau}, $\Gamma_0$ முரணுடையது என்பதைக் காட்டுகிறது.
  \end{enumerate}
\end{proof}

\end{document}
